\chapter{孤立奇点与留数}
\label{chap:isolated-singularities-residues-solutions}

\begin{solution}{laurent-geometric-annuli}
当 \(0<|z|<1\) 时，
\[
 \frac1{z(1-z)}
 =\frac1z\sum_{n=0}^{\infty}z^n
 =\sum_{n=0}^{\infty}z^{n-1}.
\]
其正向几何级数要求 \(|z|<1\)，而因子 \(1/z\) 排除中心。当 \(|z|>1\) 时，
\[
 \frac1{z(1-z)}
 =-\frac1{z^2}\frac1{1-z^{-1}}
 =-\sum_{n=0}^{\infty}z^{-n-2}.
\]
这一级数要求 \(|z^{-1}|<1\)。两个展开的公共函数相同，但环域不相交，
Laurent 系数也不同。
\end{solution}

\begin{solution}{laurent-shifted-rational}
先作部分分式
\[
 \frac1{(z-1)(z-2)}=-\frac1{z-1}+\frac1{z-2}.
\]
在 \(|z|<1\) 中，
\[
 \frac1{(z-1)(z-2)}
 =\sum_{n=0}^{\infty}\left(1-\frac1{2^{n+1}}\right)z^n.
\]
在 \(1<|z|<2\) 中，
\[
 \frac1{(z-1)(z-2)}
 =-\sum_{n=0}^{\infty}z^{-n-1}
   -\sum_{n=0}^{\infty}\frac{z^n}{2^{n+1}}.
\]
在 \(|z|>2\) 中，
\[
 \frac1{(z-1)(z-2)}
 =\sum_{n=0}^{\infty}(2^n-1)z^{-n-1}.
\]
最后一式的 \(n=0\) 系数为零。三个最大环域的边界正是到两个奇点
\(1,2\) 的距离。
\end{solution}

\begin{solution}{classify-elementary-singularities}
由幂级数
\[
 \frac{\sin z}{z^3}=\frac1{z^2}-\frac16+\frac{z^2}{120}-\cdots,
\]
故 \(0\) 是二阶极点，留数为零。又
\[
 \frac{1-\cos z}{z^4}
 =\frac1{2z^2}-\frac1{24}+\frac{z^2}{720}-\cdots,
\]
同样是二阶极点且留数为零。最后
\[
 e^{1/z}=\sum_{n=0}^{\infty}\frac1{n!z^n}
\]
有无穷多个非零负次项，所以 \(0\) 是本性奇点；\(z^{-1}\) 的系数为 \(1\)，
故留数为 \(1\)。
\end{solution}

\begin{solution}{residue-simple-quotient}
Taylor 展开给
\[
 \psi(z)=(z-a)\bigl(\psi'(a)+(z-a)h(z)\bigr),
\]
其中括号内函数在 \(a\) 处非零。于是
\[
 \frac{\phi(z)}{\psi(z)}
 =\frac1{z-a}\,
 \frac{\phi(z)}{\psi'(a)+(z-a)h(z)}.
\]
后一因子全纯，其在 \(a\) 的值为 \(\phi(a)/\psi'(a)\)。因此 Laurent
展开中 \((z-a)^{-1}\) 的系数就是该值。
\end{solution}

\begin{solution}{removable-small-order}
令 \(g(z)=(z-a)f(z)\)。题设使 \(g(z)\to0\)，故令 \(g(a)=0\) 后连续。
它在删心邻域全纯，并在 \(a\) 附近有界，所以可去奇点定理使 \(g\) 全纯。
若这一延拓恒为零，则 \(f\) 在删心邻域恒为零，直接令 \(f(a)=0\) 即得全纯
延拓。否则，\(g(a)=0\) 的孤立零点分解给
\(g(z)=(z-a)^mh(z)\)，其中 \(m\ge1\) 且 \(h(a)\ne0\)。于是
\[
 f(z)=(z-a)^{m-1}h(z)
\]
在 \(a\) 可去。若只知 \((z-a)f(z)\) 有界，取
\(f(z)=1/(z-a)\)，乘积恒为 \(1\)，而 \(a\) 是单极点，故条件不够。
\end{solution}

\begin{solution}{pole-reciprocal-zero}
若 \(f=(z-a)^{-m}g\)，其中 \(g(a)\ne0\)，则
\[
 \frac1f=(z-a)^m\frac1g
\]
延拓为全纯函数，并在 \(a\) 有恰好 \(m\) 阶零点。反之，若
\(1/f=(z-a)^mh(z)\)，其中 \(h(a)\ne0\)，则
\[
 f=(z-a)^{-m}\frac1h,
\]
\(1/h\) 在 \(a\) 附近全纯且非零，所以 \(f\) 有 \(m\) 阶极点。
\end{solution}

\begin{solution}{residue-higher-pole}
令 \(g(z)=(z-a)^mf(z)\)。它在 \(a\) 附近全纯，写
\[
 g(z)=\sum_{n=0}^{\infty}\frac{g^{(n)}(a)}{n!}(z-a)^n.
\]
除以 \((z-a)^m\) 后，\((z-a)^{-1}\) 来自 \(n=m-1\)，故
\[
 \operatorname{Res}(f,a)=\frac{g^{(m-1)}(a)}{(m-1)!}.
\]
对 \(e^z/(z-1)^4\)，有 \(m=4\)、\(g=e^z\)，所以留数为
\[
 \frac{g^{(3)}(1)}{3!}=\frac e6.
\]
\end{solution}

\begin{solution}{logarithmic-derivative-local-residue}
若 \(a\) 是 \(f\) 的 \(m\) 阶零点，写
\(f(z)=(z-a)^mg(z)\)，其中 \(g(a)\ne0\)，则
\[
 \frac{f'}f=\frac m{z-a}+\frac{g'}g.
\]
后一项在 \(a\) 全纯，所以留数为 \(m\)。若 \(a\) 是 \(p\) 阶极点，则
\(f=(z-a)^{-p}g\)，同样计算给留数 \(-p\)。若既无零点也无极点，留数为
零。因此统一写成“零点阶数减极点阶数”。
\end{solution}

\begin{solution}{residue-infinity-coefficient}
设无穷远 Laurent 展开为
\[
 f(z)=\sum_{n=-\infty}^{\infty}c_nz^n.
\]
代入 \(z=1/w\) 并除以 \(w^2\)，得到
\[
 \frac{f(1/w)}{w^2}=\sum_{n=-\infty}^{\infty}c_nw^{-n-2}.
\]
\(w^{-1}\) 项对应 \(-n-2=-1\)，即 \(n=-1\)。按定义中的负号，
\(\operatorname{Res}(f,\infty)=-c_{-1}\)。

对有理函数 \(P/Q\)，先作多项式除法，再在 \(w=1/z=0\) 对真分式作 Taylor
展开，读出 \(z^{-1}\) 系数。若真分式满足
\(\deg P\le\deg Q-2\)，则它在无穷远处是 \(O(z^{-2})\)，因而
\(c_{-1}=0\)。若 \(\deg P=\deg Q-1\)，则 \(c_{-1}\) 等于分子、分母
首项系数之比；例如
\(1/(z-a)=z^{-1}+az^{-2}+\cdots\)，故其无穷远留数为 \(-1\)。一般也可
先求全部有限极点留数，再用
\(\operatorname{Res}(f,\infty)=-\sum_{a\in\C}\operatorname{Res}(f,a)\)。
\end{solution}

\begin{solution}{real-rational-residue-integral}
上半平面只有二阶极点 \(z=i\)。写
\[
 \frac1{(z^2+1)^2}
 =\frac1{(z-i)^2(z+i)^2}.
\]
留数为
\[
 \left.\frac d{dz}(z+i)^{-2}\right|_{z=i}
 =-\frac2{(2i)^3}=-\frac i4.
\]
在半径 \(R>2\) 的上半圆上，
\[
 |z^2+1|\ge |z|^2-1\ge\frac34R^2,
\]
故被积函数模不超过 \(16/(9R^4)\)，弧积分模至多
\(16\pi/(9R^3)\to0\)。留数定理于是给
\[
 \int_{-\infty}^{\infty}\frac{dx}{(x^2+1)^2}
 =2\pi i\left(-\frac i4\right)=\frac\pi2.
\]
\end{solution}

\begin{solution}{casorati-no-limit}
若 \(f(z)\to L\in\C\)，则充分小删心邻域的像包含在
\(B(L,1)\) 中，不可能在 \(\C\) 稠密。若 \(f(z)\to\infty\)（球面意义），
则充分小删心邻域上 \(|f(z)|>1\)，其像避开圆盘 \(B(0,1)\)，也不稠密。
两种情形都与 Casorati--Weierstrass 定理矛盾。
\end{solution}

\begin{solution}{trigonometric-residue-integral}
令 \(z=e^{it}\)，则 \(dt=dz/(iz)\) 且
\(\cos t=(z+z^{-1})/2\)。积分化为
\[
 \int_{|z|=1}\frac{2\,dz}{i(bz^2+2az+b)}.
\]
令 \(\Delta=\sqrt{a^2-b^2}\)。二次式的根为
\[
 z_\pm=\frac{-a\pm\Delta}{b}.
\]
因为
\[
 |z_+|=\frac{a-\Delta}{|b|}=\frac{|b|}{a+\Delta}<1,\qquad
 |z_-|>1,
\]
单位圆内只有 \(z_+\)。被积函数在该点的留数为
\[
 \frac2{i(2bz_++2a)}
 =\frac1{i\Delta}.
\]
因此原积分为
\[
 2\pi i\frac1{i\Delta}
 =\frac{2\pi}{\sqrt{a^2-b^2}}.
\]
\end{solution}

\begin{solution}{cotangent-lattice-sum}
令
\[
 F(z)=\frac{\pi\cot(\pi z)}{z^2+a^2}.
\]
整数 \(n\) 处留数为 \(1/(n^2+a^2)\)。另外在 \(z=ia\) 处，
\[
 \operatorname{Res}(F,ia)
 =\frac{\pi\cot(\pi ia)}{2ia}
 =-\frac{\pi}{2a}\coth(\pi a),
\]
因为 \(\cot(ix)=-i\coth x\)。在 \(-ia\) 处得到同一个留数。

取边经过半整数且包住 \(\pm ia\) 的正方形轮廓。其上
\(\cot(\pi z)\) 一致有界，而 \((z^2+a^2)^{-1}=O(N^{-2})\)，故边界积分
为 \(O(N^{-1})\to0\)。留数定理及级数绝对收敛给
\[
 0=\sum_{n\in\Z}\frac1{n^2+a^2}
 -\frac{\pi}{a}\coth(\pi a),
\]
即所求公式。
\end{solution}

\begin{solution}{keyhole-beta-integral}
在 \(0<\arg z<2\pi\) 上定义
\[
 z^{\alpha-1}
 =\exp((\alpha-1)(\log|z|+i\arg z)).
\]
取沿正实轴上岸从 \(\varepsilon\) 到 \(R\)、外圆逆时针、下岸从 \(R\)
回到 \(\varepsilon\)、内圆顺时针的钥匙孔轮廓。上岸函数值为
\(x^{\alpha-1}\)，下岸为
\(e^{2\pi i(\alpha-1)}x^{\alpha-1}=e^{2\pi i\alpha}x^{\alpha-1}\)；
方向相反，所以两岸之和为
\[
 (1-e^{2\pi i\alpha})
 \int_\varepsilon^R\frac{x^{\alpha-1}}{1+x}\,dx.
\]
外圆上被积函数为 \(O(R^{\alpha-2})\)，弧长 \(2\pi R\)，故贡献
\(O(R^{\alpha-1})\to0\)。内圆上分母有正下界，贡献
\(O(\varepsilon^\alpha)\to0\)。

轮廓内只有简单极点 \(-1\)，其辐角为 \(\pi\)，留数为
\[
 (-1)^{\alpha-1}=e^{i\pi(\alpha-1)}.
\]
因此令 \(\varepsilon\downarrow0\)、\(R\to\infty\)，得到
\[
 (1-e^{2\pi i\alpha})I
 =2\pi i e^{i\pi(\alpha-1)}.
\]
利用
\[
 1-e^{2\pi i\alpha}
 =-2i e^{i\pi\alpha}\sin(\pi\alpha)
\]
化简即得 \(I=\pi/\sin(\pi\alpha)\)。
\end{solution}
